% GENERATED FILE. Edit canonical lessons, then run npm run generate:books.
% canonical-source-hash: fnv1a64:b35714871bee2748
% canonical-lessons: MW-C106-thelno, MW-C106-uthavno, MW-C106-gerno, MW-C106-balno, MW-C106-sivno

\chapter{To push, To lift, To drop, To burn, To sew}
\label{ch:mw-to-push-to-lift-to-drop-to-burn-to-sew}
\hlchaptermodality{106}

\begin{chapteropening}
\textbf{By the end of this chapter:} \emph{I can say to push, to lift, to drop, to burn and to sew.}

Complete the last lesson of chapter 106: I can say to push, to lift, to drop, to burn and to sew.
\end{chapteropening}

\section[\texorpdfstring{\mw{ठेलणो} (ṭhelṇo)}{ठेलणो (ṭhelṇo)}]{\mw{ठेलणो} (ṭhelṇo) — to push}

\label{lesson:MW-C106-thelno}



\begin{quote}
Before the new one: say the Marwadi for to fill, then the Marwadi for to pull.
\end{quote}

\subsection*{You'll want to know: \mw{ठेलणो}}
\textbf{\mw{ठेलणो}} — \emph{ṭhelṇo} — \textquotedblleft{}to push\textquotedblright{}.

\begin{cousinweb}
\textbf{\mw{ठेलणो}} (\emph{ṭhelṇo}) is a tadbhava, a word inherited and reshaped by sound change. It comes from \emph{*ṭhell}, \textquotedblleft{}push\textquotedblright{}, a form CDIAL reconstructs from its descendants. CDIAL lists Old Marwari \emph{ṭhelijaï} under it (CDIAL 5512).
\end{cousinweb}

\begin{grammarlens}[title={What you've built}]
One more word for this chapter.
\end{grammarlens}

\subsection*{Your turn}
\begin{itemize}
\raggedright
  \item \emph{Say it:} \emph{ṭhelṇo}
  \item \emph{Say it:} \emph{ṭhelṇo}, once more
  \item \emph{Say it:} say \emph{ṭhelṇo}
  \item \emph{Recall:} say the Marwadi for to fill, then the Marwadi for to pull, then say \emph{ṭhelṇo} again
\end{itemize}

\begin{tcolorbox}[breakable,colback=teal!4,colframe=teal!35!black,title={Before you move on}]
What does \mw{ठेलणो} mean? (\textquotedblleft{}To push\textquotedblright{}.) Say it once more.
\end{tcolorbox}
\section[\texorpdfstring{\mw{उठावणो} (uṭhāvṇo)}{उठावणो (uṭhāvṇo)}]{\mw{उठावणो} (uṭhāvṇo) — to lift}

\label{lesson:MW-C106-uthavno}



\begin{quote}
Before the new one: say the Marwadi for to pull, then the Marwadi for to push.
\end{quote}

\subsection*{You'll want to know: \mw{उठावणो}}
\textbf{\mw{उठावणो}} — \emph{uṭhāvṇo} — \textquotedblleft{}to lift\textquotedblright{}.

\begin{cousinweb}
\textbf{\mw{उठावणो}} (\emph{uṭhāvṇo}) is a tadbhava, a word inherited and reshaped by sound change. It comes from \emph{*ut-sthāpayati}, \textquotedblleft{}raises\textquotedblright{}, a form CDIAL reconstructs from its descendants. CDIAL lists Old Marwari \emph{uṭhāi} under it (CDIAL 1903). It is the causing form of \emph{uṭhṇo}, \textquotedblleft{}to get up\textquotedblright{}.
\end{cousinweb}

\begin{grammarlens}[title={What you've built}]
One more word for this chapter.
\end{grammarlens}

\subsection*{Your turn}
\begin{itemize}
\raggedright
  \item \emph{Say it:} \emph{uṭhāvṇo}
  \item \emph{Say it:} \emph{uṭhāvṇo}, once more
  \item \emph{Say it:} say \emph{uṭhāvṇo}
  \item \emph{Recall:} say the Marwadi for to pull, then the Marwadi for to push, then say \emph{uṭhāvṇo} again
\end{itemize}

\begin{tcolorbox}[breakable,colback=teal!4,colframe=teal!35!black,title={Before you move on}]
What does \mw{उठावणो} mean? (\textquotedblleft{}To lift\textquotedblright{}.) Say it once more.
\end{tcolorbox}
\section[\texorpdfstring{\mw{गेरणो} (gerṇo)}{गेरणो (gerṇo)}]{\mw{गेरणो} (gerṇo) — to drop}

\label{lesson:MW-C106-gerno}



\begin{quote}
Before the new one: say the Marwadi for to push, then the Marwadi for to lift.
\end{quote}

\subsection*{You'll want to know: \mw{गेरणो}}
\textbf{\mw{गेरणो}} — \emph{gerṇo} — \textquotedblleft{}to drop\textquotedblright{}.

\begin{cousinweb}
\textbf{\mw{गेरणो}} (\emph{gerṇo}) is probably a causing form, \textquotedblleft{}make fall\textquotedblright{}, of \emph{*girati}, \textquotedblleft{}drips, falls\textquotedblright{}, a form CDIAL reconstructs. CDIAL lists Old Marwari \emph{giraï}, \textquotedblleft{}falls\textquotedblright{}, under it, with Old Hindi \emph{gernā}, \textquotedblleft{}to throw\textquotedblright{}, and Gujarati \emph{geravvũ}, causing forms with the same \emph{e} as \textbf{\mw{गेरणो}} (\emph{gerṇo}) (CDIAL 4159).
\end{cousinweb}

\begin{grammarlens}[title={What you've built}]
One more word for this chapter.
\end{grammarlens}

\subsection*{Your turn}
\begin{itemize}
\raggedright
  \item \emph{Say it:} \emph{gerṇo}
  \item \emph{Say it:} \emph{gerṇo}, once more
  \item \emph{Say it:} say \emph{gerṇo}
  \item \emph{Recall:} say the Marwadi for to push, then the Marwadi for to lift, then say \emph{gerṇo} again
\end{itemize}

\begin{tcolorbox}[breakable,colback=teal!4,colframe=teal!35!black,title={Before you move on}]
What does \mw{गेरणो} mean? (\textquotedblleft{}To drop\textquotedblright{}.) Say it once more.
\end{tcolorbox}
\section[\texorpdfstring{\mw{बाळणो} (bāḷṇo)}{बाळणो (bāḷṇo)}]{\mw{बाळणो} (bāḷṇo) — to burn}

\label{lesson:MW-C106-balno}



\begin{quote}
Before the new one: say the Marwadi for to lift, then the Marwadi for to drop.
\end{quote}

\subsection*{You'll want to know: \mw{बाळणो}}
\textbf{\mw{बाळणो}} — \emph{bāḷṇo} — \textquotedblleft{}to burn\textquotedblright{}.

\begin{cousinweb}
\textbf{\mw{बाळणो}} (\emph{bāḷṇo}) is a tadbhava, a word inherited and reshaped by sound change. It comes from \emph{*dvālayati}, \textquotedblleft{}kindles\textquotedblright{}, a form CDIAL reconstructs from its descendants. CDIAL gives no Marwari form in this entry (CDIAL 6671); the match is through Hindi \emph{bālnā} and Gujarati \emph{bāḷvũ}. CDIAL sets it beside Sanskrit \emph{jvālayati}, \textquotedblleft{}kindles\textquotedblright{}. Its plain partner \emph{*dvalati}, \textquotedblleft{}burns\textquotedblright{}, has Old Marwari \emph{balaï} (CDIAL 6654).
\end{cousinweb}

\begin{grammarlens}[title={What you've built}]
One more word for this chapter.
\end{grammarlens}

\subsection*{Your turn}
\begin{itemize}
\raggedright
  \item \emph{Say it:} \emph{bāḷṇo}
  \item \emph{Say it:} \emph{bāḷṇo}, once more
  \item \emph{Say it:} say \emph{bāḷṇo}
  \item \emph{Recall:} say the Marwadi for to lift, then the Marwadi for to drop, then say \emph{bāḷṇo} again
\end{itemize}

\begin{tcolorbox}[breakable,colback=teal!4,colframe=teal!35!black,title={Before you move on}]
What does \mw{बाळणो} mean? (\textquotedblleft{}To burn\textquotedblright{}.) Say it once more.
\end{tcolorbox}
\section[\texorpdfstring{\mw{सीवणो} (sīvṇo)}{सीवणो (sīvṇo)}]{\mw{सीवणो} (sīvṇo) — to sew}

\label{lesson:MW-C106-sivno}



\begin{quote}
Before the new one: say the Marwadi for to drop, then the Marwadi for to burn.
\end{quote}

\subsection*{You'll want to know: \mw{सीवणो}}
\textbf{\mw{सीवणो}} — \emph{sīvṇo} — \textquotedblleft{}to sew\textquotedblright{}.

\begin{cousinweb}
\textbf{\mw{सीवणो}} (\emph{sīvṇo}) is a tadbhava, a word inherited from Sanskrit and reshaped by sound change. It comes from \emph{sīvyati}, \textquotedblleft{}sews\textquotedblright{}. CDIAL gives no Marwari form in this entry (CDIAL 13444); the match is through Hindi \emph{sīnā} and Gujarati \emph{sivvũ}. English \emph{sew} is the same ancient word.
\end{cousinweb}

\begin{grammarlens}[title={What you've built}]
One more word for this chapter.
\end{grammarlens}

\subsection*{Your turn}
\begin{itemize}
\raggedright
  \item \emph{Say it:} \emph{sīvṇo}
  \item \emph{Say it:} \emph{sīvṇo}, once more
  \item \emph{Say it:} say \emph{sīvṇo}
  \item \emph{Recall:} say the Marwadi for to drop, then the Marwadi for to burn, then say \emph{sīvṇo} again
\end{itemize}

\begin{tcolorbox}[breakable,colback=teal!4,colframe=teal!35!black,title={Before you move on}]
What does \mw{सीवणो} mean? (\textquotedblleft{}To sew\textquotedblright{}.) Say it once more.
\end{tcolorbox}
